\newpage
\noindent\textbf{APÊNDICE B -- NOTAS DE ENCAMINHAMENTO}
\vspace{0.5cm}

\noindent As notas deste apêndice registram a contextualização técnica, a formulação do problema e o encaminhamento metodológico que embasam a proposta de dissertação no Programa de Pós-Graduação em Engenharia de Controle e Automação (ProPECAut/Ifes).

\vspace{0.3cm}
\noindent\textbf{INTRODUÇÃO E CONTEXTUALIZAÇÃO}

\noindent A revolução proporcionada por plataformas espaciais miniaturizadas — em particular nanossatélites (CubeSats, $10 \times 10 \times 10$~cm, 1U a 3U) e picosatélites (PocketQubes, $5 \times 5 \times 5$~cm, 1p a 3p) — democratizou o acesso ao espaço para experimentos científicos, telecomunicações e sensoriamento remoto. Contudo, essa extrema miniaturização impõe restrições severas ao Subsistema de Suprimento de Energia (\textit{Electrical Power System} -- EPS). A área superficial diminuta limita a potência média gerada pelos painéis solares a valores entre centenas de miliwatts e poucos watts. Além disso, em órbita baixa da Terra (\textit{Low Earth Orbit} -- LEO), o satélite percorre cerca de 15 órbitas diárias ($\approx 90$~minutos cada), enfrentando períodos de eclipse de aproximadamente 35 a 40 minutos e variações térmicas drásticas ($-40^\circ$C a $+80^\circ$C), que alteram bruscamente o ponto de operação das células fotovoltaicas. Nesse ambiente, cada miliwatt dissipado em conversões ineficientes ou consumido por circuitos auxiliares compromete a disponibilidade operacional da carga útil (\textit{payload}) e a viabilidade da missão.

\vspace{0.3cm}
\noindent\textbf{PROBLEMA DE PESQUISA}

\noindent Como maximizar a eficiência energética global e a disponibilidade de potência no Subsistema de Energia Elétrica (EPS) de pico e nanossatélites, integrando modelagem fotovoltaica orbital, rastreamento de máxima potência (MPPT) com ultra-baixo consumo estático, conversores DC-DC de alto rendimento e algoritmos preditivos de gerenciamento dinâmico de carga?

\vspace{0.3cm}
\noindent\textbf{EIXOS METODOLÓGICOS E LINHAS DE ATUAÇÃO}

\noindent A pesquisa se desenvolverá por meio de abordagem aplicada, teórico-computacional e experimental, articulada nos seguintes eixos:
\begin{itemize}[leftmargin=*]
    \item \textbf{Eixo 1 -- Modelagem Fotovoltaica e MPPT Espacial:} Simulação computacional das curvas características I-V e P-V das células solares sob irradiação e temperatura variantes ao longo da trajetória orbital em LEO. Investigação de algoritmos de MPPT com compromisso ótimo entre acurácia de rastreamento, robustez a variações térmicas e consumo estático mínimo da eletrônica de sensoriamento e controle.
    \item \textbf{Eixo 2 -- Eletrônica de Potência e Topologias de Conversores:} Análise de perdas por condução e comutação em conversores estáticos chaveados (Buck, Boost, Buck-Boost síncronos). Projeto de topologias otimizadas para interligação entre células solares, reguladores de barramento e banco de baterias, minimizando a dissipação de calor em confinamento a vácuo.
    \item \textbf{Eixo 3 -- Gerenciamento Inteligente de Demanda e Baterias (EMS):} Desenvolvimento de lógica embarcada para despacho supervisionado de energia com base na estimativa em tempo real do Estado de Carga (SoC) e do Estado de Saúde (SoH) das células de lítio. Implementação de estratégias de \textit{duty cycling} e chaveamento dinâmico de subsistemas conforme o balanço energético em tempo real.
\end{itemize}

\vspace{0.3cm}
\noindent\textbf{CRONOGRAMA PRELIMINAR}

\begin{table}[htbp]
\centering
\small
\caption{Cronograma preliminar de execução da pesquisa (24 meses).}
\begin{tabularx}{\textwidth}{X p{3.5cm}}
\toprule
\textbf{Atividade} & \textbf{Período Previsto} \\
\midrule
Levantamento bibliográfico sistemático e análise de arquiteturas EPS & Meses 1--4 \\
Modelagem matemática orbital da matriz solar e simulação em software & Meses 3--6 \\
Projeto das topologias de conversores DC-DC e simulação eletrotérmica & Meses 5--8 \\
Implementação e validação em simulação de algoritmos de MPPT de baixo consumo & Meses 7--10 \\
Elaboração do texto de qualificação e consolidação dos modelos & Meses 10--12 \\
\textbf{Exame de Qualificação} & \textbf{Mês 12} \\
\midrule
Montagem da bancada experimental de testes e emulação solar & Meses 13--16 \\
Implementação dos algoritmos de gerenciamento de carga (EMS) em bancada & Meses 15--18 \\
Ensaios experimentais com perfis de órbita LEO e análise de rendimento global & Meses 17--20 \\
Redação de artigos científicos e escrita final da dissertação & Meses 19--23 \\
\textbf{Defesa da Dissertação} & \textbf{Mês 24} \\
\bottomrule
\end{tabularx}
\label{tab:cronograma}
\end{table}